\documentclass[a4paper,14pt]{extarticle}

\input{C:/preamble.tex}

\usepackage{tcolorbox}

\usepackage{amsmath}
\usepackage{amssymb}
\usepackage{mathtools}

\begin{document}
	Разобьём границу области $\Omega$ -- $\partial \Omega$  на твёрдые стенки $\Gamma_D$,
вход $\Gamma_{\mathrm{L}}$ и выход $\Gamma_{\mathrm{R}}$.


Пусть $\mathbf n = (n_x,n_y)$ -- внешняя единичная нормаль
к $\partial \Omega$ . Введём функции
$
	p_{\mathrm{L}}(y),~ p_{\mathrm{R}}(y)
$
на $\Gamma_{\mathrm{L}}$ и $\text{на } \Gamma_{\mathrm{R}}$.







Тогда стационарная задача Стокса для постоянной вязкости
с условием прилипания на стенках и заданным псевдонапряжением
на входе и выходе имеет вид
\begin{equation}
	\begin{cases}
		-\nu \Delta v_x
		+ \dfrac{\partial p}{\partial x}
		= f_x, \\

		-\nu \Delta v_y
		+ \dfrac{\partial p}{\partial y}
		= f_y, \\

		\dfrac{\partial v_x}{\partial x}
		+
		\dfrac{\partial v_y}{\partial y}
		= 0,
	\end{cases}
	\label{eq1}
	\end{equation}
	с граничными условиями
	\begin{equation}
		\begin{gathered}
		v_x \big |_{\Gamma_D} = 0,\quad
		v_y \big |_{\Gamma_D} = 0,\\
		\nu\dfrac{\partial v_x}{\partial n}
		- p n_x
		= -p_{\mathrm{L}} n_x, \quad
		\nu\dfrac{\partial v_x}{\partial n}
		- p n_x
		= -p_{\mathrm{R}} n_x, \\
		\nu\dfrac{\partial v_y}{\partial n}
		- p n_y
		= -p_{\mathrm{L}} n_y, \quad
		\nu\dfrac{\partial v_y}{\partial n}
		- p n_y
		= -p_{\mathrm{R}} n_y,
	\end{gathered}
	\end{equation}
	здесь $v_x$ и $v_y$ --- компоненты скорости,
	$p$ --- давление, $\nu$ --- известная постоянная вязкость,
	а $f_x$ и $f_y$ --- компоненты заданной объёмной силы.
	
	Принимая
	\begin{equation}
	{\mathbf u} = (v_x,v_y)^{T},\quad
	\mathbf f = (f_x, f_y)^{T},\quad
	\Gamma_{N} = \Gamma_{\mathrm{L}} \cup \Gamma_{\mathrm{R}}, \quad
	p_{b} = p_{\mathrm{L}} \cup p_{\mathrm{R}},
	\end{equation}
	оно записывается в векторном виде
	\begin{equation}
		\begin{dcases}
			-\nu\Delta{\mathbf u}
			+\nabla p
			=\mathbf f,
			\\
			\operatorname{div}{\mathbf u}=0,
		\end{dcases}
		\qquad
		\begin{gathered}
			\left.{\mathbf u}\right|_{\Gamma_D}=0,
			\\
			\left.
			\left(
			\nu\frac{\partial{\mathbf u}}{\partial n}
			-p\mathbf n
			\right)
			\right|_{\Gamma_N}
			=
			-p_b\mathbf n.
		\end{gathered}
		\label{eq:stokes_vector}
	\end{equation}
	
	
	\section{Слабая форма}
	
	Домножим первое и второе уравнения~\eqref{eq1}
	на тестовые функции $\varphi$ и $\psi$, которые равны 0 на
	$\Gamma_D$ (на $\Gamma_N$ они могут быть ненулевыми), а третье уравнение --- на $q$, и проинтегрируем по всей области.

	
	Первое уравнение примет  вид
	\begin{equation}
		-\nu \int_{\Omega}
		\left(
		\frac{\partial^2 v_x}{\partial x^2}
		+
		\frac{\partial^2 v_x}{\partial y^2}
		\right)
		\varphi \, d\Omega
		+
		\int_{\Omega}
		\frac{\partial p}{\partial x}
		\varphi \, d\Omega
		=
		\int_{\Omega} f_x\varphi \, d\Omega.
	\end{equation}
	
	Проинтегрируем слагаемые левой части по частям:
	\begin{enumerate}
		\item Для слагаемого с лапласианом:
		\begin{equation*}
			\int_{\Omega}
			\Delta v_x\,\varphi \, d\Omega
			=
			-\int_{\Omega}
			\nabla v_x\cdot\nabla\varphi \, d\Omega
			+
			\int_{\partial\Omega}
			\frac{\partial v_x}{\partial n}
			\varphi \, d\Gamma.
		\end{equation*}
		
		\item Для слагаемого с давлением:
		\begin{equation*}
			\int_{\Omega}
			\frac{\partial p}{\partial x}
			\varphi \, d\Omega
			=
			-\int_{\Omega}
			p\frac{\partial\varphi}{\partial x} \, d\Omega
			+
			\int_{\partial\Omega}
			p\,n_x\varphi \, d\Gamma.
		\end{equation*}
	\end{enumerate}
	
	Поскольку $\varphi=0$ на $\Gamma_D$, в граничных интегралах
	остаётся только участок $\Gamma_N$. Объединяя граничные
	слагаемые и перенося их в правую часть, получаем
	\begin{equation*}
		\begin{aligned}
			\nu\int_{\Omega}
			\nabla v_x\cdot\nabla\varphi \, d\Omega
			-
			\int_{\Omega}
			p\frac{\partial\varphi}{\partial x} \, d\Omega
			={}&
			\int_{\Omega} f_x\varphi \, d\Omega
			\\
			&+
			\int_{\Gamma_N}
			\left(
			\nu\frac{\partial v_x}{\partial n}
			-p n_x
			\right)
			\varphi \, d\Gamma.
		\end{aligned}
	\end{equation*}
	
	Подставляя граничное условие и аналогично преобразуя второе
	уравнение, получаем слабую форму:
	\begin{equation}
		\begin{dcases}
			\nu\int_\Omega
			\nabla v_x\cdot\nabla\varphi \, d\Omega
			-
			\int_\Omega
			p\frac{\partial\varphi}{\partial x} \, d\Omega
			=
			\int_\Omega f_x\varphi \, d\Omega
			-
			\int_{\Gamma_N}
			p_b n_x\varphi \, d\Gamma,
			\\
			\nu\int_\Omega
			\nabla v_y\cdot\nabla\psi \, d\Omega
			-
			\int_\Omega
			p\frac{\partial\psi}{\partial y} \, d\Omega
			=
			\int_\Omega f_y\psi \, d\Omega
			-
			\int_{\Gamma_N}
			p_b n_y\psi \, d\Gamma,
			\\
			\int_\Omega q\operatorname{div}u \, d\Omega = 0,
		\end{dcases}
		\label{eq:weak}
	\end{equation}
	где $v_x,v_y\in H^1(\Omega)$, $p\in L^2(\Omega)$,
	причём $v_x=v_y=0$ на $\Gamma_D$. Равенства должны
	выполняться для всех $\varphi,\psi\in H^1(\Omega)$,
	равных нулю на $\Gamma_D$, и всех $q\in L^2(\Omega)$.
	Граничные равенства понимаются в смысле следа.

\section{Разбиение области и тестовые функции}

Разобьём область $\Omega$ на $N_\tau$ подобластей
$\tau_\alpha$ с непересекающимися внутренностями:
\begin{equation}
	\overline{\Omega}
	=
	\bigcup_{\alpha=1}^{N_\tau}
	\overline{\tau_\alpha}.
\end{equation}

Обозначим через $\gamma_\alpha$ часть границы Неймана,
принадлежащую границе подобласти $\tau_\alpha$:
\begin{equation}
	\gamma_\alpha
	=
	\partial\tau_\alpha \cap \Gamma_N.
\end{equation}

Предположим, что участки $\gamma_\alpha$ покрывают $\Gamma_N$
и не перекрываются по длине. Если подобласть $\tau_\alpha$
не соприкасается с $\Gamma_N$ по участку ненулевой длины,
то соответствующий граничный вклад равен нулю.

Интеграл по области представим суммой интегралов по подобластям,
а граничный интеграл --- суммой интегралов по участкам
$\gamma_\alpha$. Применяя это к первому уравнению слабой
формы~\eqref{eq:weak}, получаем
\begin{equation}
	\begin{gathered}
		\nu\sum_{\alpha=1}^{N_\tau}\int_{\tau_\alpha}
		\nabla v_x\cdot\nabla\varphi\,d\tau_\alpha
		-
		\sum_{\alpha=1}^{N_\tau}\int_{\tau_\alpha}
		p\frac{\partial\varphi}{\partial x}\,d\tau_\alpha
		=\\
		=
		\sum_{\alpha=1}^{N_\tau}\int_{\tau_\alpha}
		f_x\varphi\,d\tau_\alpha
		-
		\sum_{\alpha=1}^{N_\tau}\int_{\gamma_\alpha}
		p_b n_x\varphi\,d\Gamma.
	\end{gathered}
	\label{eq:parts_first}
\end{equation}

Равенство выполняется для каждой допустимой глобальной
тестовой функции $\varphi$. Граничное слагаемое не исчезает,
так как тестовая функция на $\Gamma_N$ может быть ненулевой.

Выберем конечные наборы допустимых тестовых функций:
\[
	\varphi_1,\ldots,\varphi_{m_\varphi},\qquad
	\psi_1,\ldots,\psi_{m_\psi},\qquad
	q_1,\ldots,q_{m_q}.
\]

Функции $\varphi_i$ и $\psi_i$ принадлежат $H^1(\Omega)$
и равны нулю на $\Gamma_D$, а $q_i\in L^2(\Omega)$.
Подставляя их в соответствующие уравнения, получаем систему
с конечным количеством строк:
\begin{equation}
	\begin{dcases}
		\nu\sum_{\alpha=1}^{N_\tau}\int_{\tau_\alpha}
		\nabla v_x\cdot\nabla\varphi_i\,d\tau_\alpha
		-
		\sum_{\alpha=1}^{N_\tau}\int_{\tau_\alpha}
		p\frac{\partial\varphi_i}{\partial x}\,d\tau_\alpha
		=\\
		=
		\sum_{\alpha=1}^{N_\tau}\int_{\tau_\alpha}
		f_x\varphi_i\,d\tau_\alpha
		-
		\sum_{\alpha=1}^{N_\tau}\int_{\gamma_\alpha}
		p_b n_x\varphi_i\,d\Gamma,
		\quad i=1,\ldots,m_\varphi;
		\\
		\cdots\quad
		\text{(второе уравнение аналогично)}
		\quad i=1,\ldots,m_\psi;
		\\
		\sum_{\alpha=1}^{N_\tau}\int_{\tau_\alpha}
		q_i\operatorname{div}\mathbf u\,d\tau_\alpha=0,
		\qquad i=1,\ldots,m_q.
	\end{dcases}
	\label{eq:parts_div}
\end{equation}

Получены равенства для выбранных тестовых функций. Искомые
функции здесь ещё не заменены конечномерными приближениями,
поэтому это пока не система относительно конечного числа
числовых коэффициентов.

Обозначим ограничения тестовых функций на подобласти:
\[
	\Phi_{i,\alpha}
	=
	\varphi_i|_{\tau_\alpha},\qquad
	\Psi_{i,\alpha}
	=
	\psi_i|_{\tau_\alpha},\qquad
	Q_{i,\alpha}
	=
	q_i|_{\tau_\alpha}.
\]

Для первого набора ограничения можно расположить в матрицу
\begin{equation}
	\Phi
	=
	\begin{pmatrix}
		\Phi_{1,1}
		&
		\cdots
		&
		\Phi_{1,N_\tau}
		\\
		\vdots
		&
		\ddots
		&
		\vdots
		\\
		\Phi_{m_\varphi,1}
		&
		\cdots
		&
		\Phi_{m_\varphi,N_\tau}
	\end{pmatrix}.
\end{equation}

Аналогично задаются матрицы ограничений $\Psi$ и $\mathcal Q$.
Каждая строка содержит ограничения одной глобальной функции;
ограничения $\Phi_{i,\alpha}$ и $\Phi_{i,\beta}$ на общей
границе согласованы в смысле следов. Аналогичное замечание
относится к $\Psi$. Для функций $q_i\in L^2(\Omega)$
совпадение следов без дополнительных требований не предполагается.

Интеграл по подобласти $\tau_\alpha$ зависит только от
значений функций внутри этой подобласти. Поэтому вместо
функций определенных на всей области можно использовать их ограничения на $\tau_\alpha$: значения функций и их производных внутри подобласти при этом не изменяются.

Тогда первое уравнение принимает вид
\begin{equation}
	\begin{gathered}
		\nu\sum_{\alpha=1}^{N_\tau}\int_{\tau_\alpha}
		\nabla\left(v_x|_{\tau_\alpha}\right)
		\cdot\nabla\Phi_{i,\alpha}\,d\tau_\alpha
		-
		\sum_{\alpha=1}^{N_\tau}\int_{\tau_\alpha}
		\left(p|_{\tau_\alpha}\right)
		\frac{\partial\Phi_{i,\alpha}}{\partial x}
		\,d\tau_\alpha
		=\\
		=
		\sum_{\alpha=1}^{N_\tau}\int_{\tau_\alpha}
		\left(f_x|_{\tau_\alpha}\right)
		\Phi_{i,\alpha}\,d\tau_\alpha
		-
		\sum_{\alpha=1}^{N_\tau}\int_{\gamma_\alpha}
		p_b n_x\Phi_{i,\alpha}\,d\Gamma,
		\quad i=1,\ldots,m_\varphi.
	\end{gathered}
	\label{eq:restriction_x}
\end{equation}



\section{Аппроксимация решения}

В рассматриваемых пространствах можно выбрать полные
ортонормированные системы функций. Согласно теореме о базисе
сепарабельного гильбертова пространства, каждую функцию
такого пространства можно представить рядом по соответствующей
системе, сходящимся в норме этого пространства
(см.~\cite{hilbert_basis}).

Следовательно, отбрасывая все слагаемые ряда после некоторого
номера, можно получить конечную линейную комбинацию,
приближающую исходную функцию. При увеличении числа слагаемых
погрешность такого приближения стремится к нулю в норме
соответствующего пространства.

Поэтому численные решения будем искать в виде конечных линейных комбинаций базисных функций.
Для скорости будем выбирать базисные функции из
$H^1(\Omega)$ с нулевым следом на $\Gamma_D$, а для давления
--- из $L^2(\Omega)$. Поэтому конечные линейные комбинации
этих функций также принадлежат соответствующим пространствам,
а приближённая скорость удовлетворяет однородному граничному
условию на $\Gamma_D$.

\begin{equation}
	v_x \approx
	v_x^*
	=
	\sum_{j=1}^{N_x}\omega_{x,j}\lambda_{x,j},
	\quad
	v_y\approx
	v_y^*
	=
	\sum_{j=1}^{N_y}\omega_{y,j}\lambda_{y,j},
	\quad
	p \approx
	p^*
	=
	\sum_{k=1}^{N_p}\omega_{p,k}\lambda_{p,k},
\end{equation}
здесь $\omega_{x,j}$, $\omega_{y,j}$ и $\omega_{p,k}$ ---
неизвестные числовые коэффициенты.
Ограничивая эти разложения на подобласть $\tau_\alpha$,
получаем
\begin{equation}
	v_x^*|_{\tau_\alpha}
	=
	\sum_{j=1}^{m_\varphi}
	\omega_{x,j}\Phi_{j,\alpha},
	\quad
	v_y^*|_{\tau_\alpha}
	=
	\sum_{j=1}^{m_\psi}
	\omega_{y,j}\Psi_{j,\alpha},
	\quad
	p^*|_{\tau_\alpha}
	=
	\sum_{k=1}^{m_q}
	\omega_{p,k}Q_{k,\alpha}.
	\label{eq:local_approximation}
\end{equation}

Теперь подставим это рахложжение, получим систему
\begin{equation}
	\begin{gathered}
		\nu\sum_{\alpha=1}^{N_\tau}
		\int_{\tau_\alpha}
		\nabla\left(
			\sum_{j=1}^{N_x}
			\omega_{x,j}\lambda_{x,j,\alpha}
		\right)
		\cdot\nabla\Phi_{i,\alpha}\,d\tau_\alpha
		-
		\sum_{\alpha=1}^{N_\tau}
		\int_{\tau_\alpha}
		\left(
			\sum_{k=1}^{N_p}
			\omega_{p,k}\lambda_{p,k,\alpha}
		\right)
		\frac{\partial\Phi_{i,\alpha}}{\partial x}
		\,d\tau_\alpha
		=\\
		=
		\sum_{\alpha=1}^{N_\tau}
		\int_{\tau_\alpha}
		f_x\Phi_{i,\alpha}\,d\tau_\alpha
		-
		\sum_{\alpha=1}^{N_\tau}
		\int_{\gamma_\alpha}
		p_b n_x\Phi_{i,\alpha}\,d\Gamma,
		\quad i=1,\ldots,m_\varphi.
	\end{gathered}
\end{equation}
Коэффициент $\omega_{x,j}$ соответствует глобальной базисной
функции $\lambda_{x,j}$ и поэтому остаётся одним и тем же
при рассмотрении её ограничений на разные подобласти.
В частности, если функция $\lambda_{x,j}$ ненулевая на
нескольких соседних подобластях, то во всех этих подобластях
при её ограничениях стоит один и тот же коэффициент
$\omega_{x,j}$. Аналогично это относится к коэффициентам
$\omega_{y,j}$ и $\omega_{p,k}$.

Поскольку коэффициенты $\omega_{x,j}$ и $\omega_{p,k}$
не зависят от координат, их можно вынести за знаки
интегралов. Используя линейность градиента и интеграла,
а также меняя порядок конечных сумм, получаем
\begin{equation}
	\begin{gathered}
		\sum_{j=1}^{N_x}
		\omega_{x,j}
		\left(
			\nu\sum_{\alpha=1}^{N_\tau}
			\int_{\tau_\alpha}
			\nabla\lambda_{x,j,\alpha}
			\cdot\nabla\Phi_{i,\alpha}
			\,d\tau_\alpha
		\right)
		-
		\sum_{k=1}^{N_p}
		\omega_{p,k}
		\left(
			\sum_{\alpha=1}^{N_\tau}
			\int_{\tau_\alpha}
			\lambda_{p,k,\alpha}
			\frac{\partial\Phi_{i,\alpha}}{\partial x}
			\,d\tau_\alpha
		\right)
		=
		\\
		=
		\sum_{\alpha=1}^{N_\tau}
		\int_{\tau_\alpha}
		f_x\Phi_{i,\alpha}\,d\tau_\alpha
		-
		\sum_{\alpha=1}^{N_\tau}
		\int_{\gamma_\alpha}
		p_b n_x\Phi_{i,\alpha}\,d\Gamma,
		\quad i=1,\ldots,m_\varphi.
	\end{gathered}
\end{equation}
Множители коэффициентов нам известны, поэтому интегралы в полученном
уравнении можно вычислить; неизвестными остаются
коэффициенты.

Введём обозначения
\begin{equation}
	\begin{gathered}
		(A_x)_{ij}
		=
		\nu\sum_{\alpha=1}^{N_\tau}
		\int_{\tau_\alpha}
		\nabla\lambda_{x,j,\alpha}
		\cdot\nabla\Phi_{i,\alpha}\,d\tau_\alpha,
		\\
		(B_x)_{ik}
		=
		\sum_{\alpha=1}^{N_\tau}
		\int_{\tau_\alpha}
		\lambda_{p,k,\alpha}
		\frac{\partial\Phi_{i,\alpha}}{\partial x}
		\,d\tau_\alpha,
		\\
		(F_x)_i
		=
		\sum_{\alpha=1}^{N_\tau}
		\int_{\tau_\alpha}
		f_x\Phi_{i,\alpha}\,d\tau_\alpha
		-
		\sum_{\alpha=1}^{N_\tau}
		\int_{\gamma_\alpha}
		p_b n_x\Phi_{i,\alpha}\,d\Gamma,\\
		\mathcal{W}_{x} = (w_{x,1},\dots,w_{x,N_x})^T,\quad
		\mathcal{W}_{p} = (w_{p,1},\dots,w_{p,N_p})^T
	\end{gathered}
\end{equation}
Тогда получается система с $N_x + N_p$ столбцами и $m_\varphi$ строками
\begin{equation}
	A_x\mathcal{W}_x - B_x\mathcal{W}_p   = F_x.
\end{equation}


Введя $A_y$, $B_y$, $F_y$, $W_y$,  получим то же самое и для второго уравнения
\begin{equation}
	A_y\mathcal W_y-B_y\mathcal W_p=F_y.
\end{equation}

Для третьего уравнения ($i=1,\ldots,m_q$)
\begin{equation}
	\begin{gathered}
		\sum_{j=1}^{N_x}
		\omega_{x,j}
		\left(
			\sum_{\alpha=1}^{N_\tau}
			\int_{\tau_\alpha}
			Q_{i,\alpha}
			\frac{\partial\lambda_{x,j,\alpha}}{\partial x}
			\,d\tau_\alpha
		\right)
		+
		\sum_{j=1}^{N_y}
		\omega_{y,j}
		\left(
			\sum_{\alpha=1}^{N_\tau}
			\int_{\tau_\alpha}
			Q_{i,\alpha}
			\frac{\partial\lambda_{y,j,\alpha}}{\partial y}
			\,d\tau_\alpha
		\right)
		=0
	\end{gathered}
\end{equation}
вводим
\begin{equation}
		(C_x)_{ij}
		=
		\sum_{\alpha=1}^{N_\tau}
		\int_{\tau_\alpha}
		Q_{i,\alpha}
		\frac{\partial\lambda_{x,j,\alpha}}{\partial x}
		\,d\tau_\alpha,\quad
		(C_y)_{ij}
		=
		\sum_{\alpha=1}^{N_\tau}
		\int_{\tau_\alpha}
		Q_{i,\alpha}
		\frac{\partial\lambda_{y,j,\alpha}}{\partial y}
		\,d\tau_\alpha.
\end{equation}
тогда
\begin{equation}
	C_x\mathcal W_x+C_y\mathcal W_y=0.
\end{equation}

Поную систему можно записать блочной матрией
\begin{equation}
	\begin{pmatrix}
		A_x & 0 & -B_x\\
		0 & A_y & -B_y\\
		C_x & C_y & 0
		\end{pmatrix}
		\begin{pmatrix}
		\mathcal W_x\\
		\mathcal W_y\\
		\mathcal W_p
		\end{pmatrix}
		=
		\begin{pmatrix}
		F_x\\
		F_y\\
		0.
	\end{pmatrix}
\end{equation}

Условия Неймана уже учтены в правых частях $F_x$ и $F_y$.
Остаётся учесть условия Дирихле: на $\Gamma_D$ обе
компоненты скорости равны нулю.

Для каждого узла скорости, расположенного на $\Gamma_D$,
соответствующие значения коэффициентов $\omega_{x,j}$ и
$\omega_{y,j}$ принимаются равными нулю. Если система
собрана для всех узлов, то в ней для каждого такого
коэффициента обнуляются соответствующие строка и столбец,
на диагонали устанавливается единица, а соответствующий
элемент правой части приравнивается к нулю.

После этого полученная система готова к решению.




\section{Выбор базисных функций}

Для построения приближённого решения воспользуемся методом
Галёркина. Выберем тестовые функции совпадающими с базисными
функциями, используемыми для приближения соответствующих
искомых функций:
$
	\varphi_i=\lambda_{x,i},
$ 
$
\psi_i=\lambda_{y,i},
$
$
	q_i=\lambda_{p,i}.
$
Тогда количество тестовых функций совпадает с количеством
соответствующих базисных функций:
$
	m_\varphi=N_x,
$
$
	m_\psi=N_y,
$
$
	m_q=N_p.
$

Тогда возникают соотношения
\begin{equation}
		(C_x)_{ki}
		=
		\sum_{\alpha=1}^{N_\tau}
		\int_{\tau_\alpha}
		\lambda_{p,k,\alpha}
		\frac{\partial\lambda_{x,i,\alpha}}{\partial x}
		\,d\tau_\alpha
		=
		(B_x)_{ik},
		\quad
		\text{и}
		\quad
		(C_y)_{ki} = (B_x)_{ik}.
\end{equation}

Таким образом, глобальную систему можно записать в виде
\begin{equation}
	\begin{pmatrix}
		A_x & 0 & -B_x\\
		0 & A_y & -B_y\\
		B_x^T & B_y^T & 0
	\end{pmatrix}
	\begin{pmatrix}
		\mathcal W_x\\
		\mathcal W_y\\
		\mathcal W_p
	\end{pmatrix}
	=
	\begin{pmatrix}
		F_x\\
		F_y\\
		0
	\end{pmatrix}.
\end{equation}
Полученная матрица является квадратной и имеет размер
$(N_x+N_y+N_p)\times(N_x+N_y+N_p)$.

Для построения устойчивой дискретной системы необходимо
подходящим образом выбрать пространства приближений скорости
и давления. Одним из таких выборов является элемент
Тейлора--Худа: область разюивается на треугольники, на каждом скорость приближается полиномами второй степени, а давление --- полиномами первой
степени. Устойчивость соответствующей пары пространств,
в том числе выполнение дискретного условия inf-sup при
определённых условиях на сетку и граничные условия,
доказана в литературе~\cite{taylor_hood}.

Имеем $N_T$ треугольников $T_h$. Поскольку
$\left.p^*\right|_{T_h}\in P_1(T_h)$ для каждого
$h=1,\ldots,N_T$, то это ограничение однозначно задаётся
тремя коэффициентами $(c_1,c_2,c_3)$ в базисе $\{1,x,y\}$.
Следовательно, для задания базиса в $P_1(T_h)$ необходимы
три линейно независимые функции.

В качестве базиса можно использовать не только $\{1,x,y\}$,
но и другие три линейно независимые функции из $P_1(T_h)$.
Выберем базис, связанный с вершинами треугольника.


Для того чтобы ограничения $p^*|_{T_a}$ и $p^*|_{T_b}$
совпадали на общем ребре $T_a\cap T_b$, достаточно
совпадения их значений в двух различных точках этого ребра,
поскольку оба ограничения линейны.

Удобно выбрать базисные функции так, чтобы коэффициенты
разложения на треугольнике $\omega_{p,i}^{(h)}$ были равны
значениям приближённого давления в его вершинах:
\begin{equation}
	p^*\left(a_j^{(h)}\right)
	=
	\omega_{p,j}^{(h)},
	\qquad j=1,2,3;
\end{equation}
подставляя разложение давления, получаем
\begin{equation}
	\sum_{i=1}^{3}\omega_{p,i}^{(h)}\lambda_{p,i}^{(h)}(a_j^{(h)})
	=\omega_{p,j}^{(h)};
\end{equation}
	отсюда получаем требования
\begin{equation}
		\lambda_{p,i}^{(h)}(a_j^{(h)})=\delta_{ij},\quad \text{и} \quad
		\lambda_{p,1}^{(h)}+	\lambda_{p,2}^{(h)}+	\lambda_{p,3}^{(h)} = 1.
\end{equation}

\begin{tcolorbox}
		Рассмотрим эталонный треугольник
		$\widehat{T}\subset\mathbb{R}^2_{\xi,\eta}$
		с вершинами
		$M_1=(0,0)$, $M_2=(1,0)$, $M_3=(0,1)$.
		На нем базисные функции
		пространства $P_1(\widehat{T})$,
		удовлетворяющие условиям~\eqref{eq:qwe},
		имеют вид
		\begin{equation}
			\phi_1(\xi,\eta)=1-\xi-\eta,\qquad
			\phi_2(\xi,\eta)=\xi,\qquad
			\phi_3(\xi,\eta)=\eta.
		\end{equation}
	\end{tcolorbox}
	
	Для построения базисных функций на произвольном
	треугольнике $T_h$ сетки можно воспользоваться
	следующим подходом:
	\begin{enumerate}
		\item ввести  $\widehat T$;
		
		\item построить обратимое аффинное
		отображение для треугльника с вершинами $(x_1,y_1),~(x_2,y_2),~(x_3,y_3)$ -- 
		$G_h:T_h\to\widehat T$;
		
		\item на $\widehat T$ задать $\phi_1,\phi_2,\phi_3$;
		
		\item с помощью обратного отображения
		$G_h^{-1}$ перенести $\phi_1,\phi_2,\phi_3$
		на  $T_h$,
		получив 
		$\lambda_{p,i}^{(h)}(x,y)$;
	\end{enumerate}
	то есть
	\begin{equation}
		\lambda^{(h)}_{p,i}(x,y) = \phi_i(G_h(x,y)),\quad x,y \in T_h.
	\end{equation}
	В результате они будут выглядеть следующим образом
	\begin{equation}
	\begin{gathered}
	D^{(h)} = (x_2 - x_1)(y_3 - y_1) - (x_3 - x_1)(y_2 - y_1),
	\\
	\lambda_{p,1}^{(h)}(x,y) =
	\frac{(x_2 - x)(y_3 - y) - (x_3 - x)(y_2 - y)}{D^{(h)}},
	\\
	\lambda_{p,2}^{(h)}(x,y) =
	\frac{(y_3 - y_1)(x - x_1) - (x_3 - x_1)(y - y_1)}{D^{(h)}},
	\\
	\lambda_{p,3}^{(h)}(x,y) =
	\frac{(x_2 - x_1)(y - y_1) - (y_2 - y_1)(x - x_1)}{D^{(h)}},
	\end{gathered}
	\quad (x,y) \in T_h.
	\end{equation}
	Коэффициенты, после приведения подобных, определяются
	\begin{equation}
		\begin{pmatrix}
			c_{1,1} & c_{1,2} & c_{1,3}\\
			c_{2,1} & c_{2,2} & c_{2,3}\\
			c_{3,1} & c_{3,2} & c_{3,3}
		\end{pmatrix}
		=
		\frac{1}{D^{(h)}}
		\begin{pmatrix}
			x_2y_3-x_3y_2 & y_2-y_3 & x_3-x_2\\
			x_3y_1-x_1y_3 & y_3-y_1 & x_1-x_3\\
			x_1y_2-x_2y_1 & y_1-y_2 & x_2-x_1
		\end{pmatrix}.
	\end{equation}

	Для приближения каждой компоненты скорости на треугольнике
	$T_h$ используется пространство $P_2(T_h)$. Его базис
	$\{1,x,y,x^2,xy,y^2\}$ содержит шесть функций, поэтому для
	построения узлового базиса выберем шесть узлов: три вершины
	треугольника и середины трёх его рёбер. Точек в которых значения (значения скоростей) на общем ребре дожны совпадать -- 3, возьмем за них вершины и середину ребера. Получаем те же условия
	\begin{equation}
	\lambda_{u,i}^{(h)}\left(a_j^{(h)}\right)
	=
	\delta_{ij},
	\quad i,j=1,\ldots,6 \quad \text{и} \quad
	\lambda_1 + \cdots + \lambda_6 = 1.
	\end{equation}
	Подставляя подставляя координаты шести узлов в это выражение, получаем систему относительно коэффициентов
	$c_{i,1},\ldots,c_{i,6}$:
	\begin{equation}
		c_{i,1}+c_{i,2}x_j+c_{i,3}y_j
		+c_{i,4}x_j^2+c_{i,5}x_jy_j+c_{i,6}y_j^2
		=
		\delta_{ij},
		\qquad j=1,\ldots,6.
	\end{equation}
	Для невырожденного треугольника эта система решается однозначно
\begin{equation}
	\begin{pmatrix}
		c_{1,1} & c_{1,2} & c_{1,3} & c_{1,4} & c_{1,5} & c_{1,6}\\
		c_{2,1} & c_{2,2} & c_{2,3} & c_{2,4} & c_{2,5} & c_{2,6}\\
		c_{3,1} & c_{3,2} & c_{3,3} & c_{3,4} & c_{3,5} & c_{3,6}\\
		c_{4,1} & c_{4,2} & c_{4,3} & c_{4,4} & c_{4,5} & c_{4,6}\\
		c_{5,1} & c_{5,2} & c_{5,3} & c_{5,4} & c_{5,5} & c_{5,6}\\
		c_{6,1} & c_{6,2} & c_{6,3} & c_{6,4} & c_{6,5} & c_{6,6}
	\end{pmatrix}
	=
	\begin{pmatrix}
		2a_1^2-a_1 & (4a_1-1)b_1 & (4a_1-1)d_1
		& 2b_1^2 & 4b_1d_1 & 2d_1^2\\

		2a_2^2-a_2 & (4a_2-1)b_2 & (4a_2-1)d_2
		& 2b_2^2 & 4b_2d_2 & 2d_2^2\\

		2a_3^2-a_3 & (4a_3-1)b_3 & (4a_3-1)d_3
		& 2b_3^2 & 4b_3d_3 & 2d_3^2\\

		4a_1a_2 & 4(a_1b_2+a_2b_1) & 4(a_1d_2+a_2d_1)
		& 4b_1b_2 & 4(b_1d_2+b_2d_1) & 4d_1d_2\\

		4a_2a_3 & 4(a_2b_3+a_3b_2) & 4(a_2d_3+a_3d_2)
		& 4b_2b_3 & 4(b_2d_3+b_3d_2) & 4d_2d_3\\

		4a_3a_1 & 4(a_3b_1+a_1b_3) & 4(a_3d_1+a_1d_3)
		& 4b_3b_1 & 4(b_3d_1+b_1d_3) & 4d_3d_1
	\end{pmatrix}.
\end{equation}







\end{document}
